\documentclass[10pt,aspectratio=43]{beamer}

% ============================================================
% Theme and style — match the Reference_Slides aesthetic
% ============================================================
\usetheme{default}
\usecolortheme{default}
\setbeamertemplate{navigation symbols}{}
\setbeamertemplate{footline}[frame number]

\usepackage{lmodern}
\usepackage[T1]{fontenc}
\usepackage{amsmath,amssymb,amsfonts}
\usepackage{booktabs}
\usepackage{array}
\usepackage{graphicx}
\usepackage{tikz}
\usetikzlibrary{shapes,arrows,positioning,calc,decorations.pathreplacing}

\usepackage{tcolorbox}
\tcbuselibrary{skins}

\newtcolorbox{yellowbox}{
  colback=yellow!20, colframe=yellow!20, boxrule=0pt, arc=0pt,
  left=4pt, right=4pt, top=4pt, bottom=4pt}

\newtcolorbox{redbox}{
  colback=red!15, colframe=red!15, boxrule=0pt, arc=0pt,
  left=4pt, right=4pt, top=4pt, bottom=4pt}

\newtcolorbox{greenbox}{
  colback=green!10, colframe=green!10, boxrule=0pt, arc=0pt,
  left=4pt, right=4pt, top=4pt, bottom=4pt}

\setbeamertemplate{footline}{%
  \hfill\insertframenumber/\inserttotalframenumber\hspace*{6pt}\vspace*{4pt}%
}
\setbeamertemplate{itemize items}[circle]

% Page reference command
\newcommand{\pref}[1]{\vfill\hfill{\scriptsize\color{gray}[Norton, pp.~#1]}}

% ============================================================
\title{Chapter 2: Kinematics Fundamentals}
\subtitle{Design of Machinery --- Norton, 6th Edition}
\author{}
\date{}

\begin{document}

% ============================================================
% TITLE
% ============================================================

\begin{frame}
\frametitle{Chapter 2: Kinematics Fundamentals}

\begin{itemize}
  \item Degrees of freedom and mobility
  \item Types of motion
  \item Links, joints, and kinematic chains
  \item Gruebler's and Kutzbach's equations
  \item Mechanisms vs.\ structures
  \item Number synthesis
  \item Paradoxes and isomers
  \item Linkage transformation
  \item Intermittent motion and inversion
  \item The Grashof condition and fourbar classification
  \item Linkages of more than four bars
  \item Practical considerations and motors
\end{itemize}

\vspace{4pt}
\textbf{Related Reading:} Norton, Ch.~2, pp.~30--97.
\end{frame}

% ============================================================
% 2.0--2.1 DOF
% ============================================================

\begin{frame}
\frametitle{Degrees of Freedom (DOF) or Mobility}

\begin{yellowbox}
\textbf{Degree of freedom (DOF)} or \textbf{mobility} ($M$) of a mechanical system: \emph{the number of independent parameters (measurements) needed to uniquely define its position in space at any instant of time.}
\end{yellowbox}

\vspace{6pt}
\textbf{A rigid body in a plane} (2-D):
\begin{itemize}
  \item Needs 3 parameters: two linear coordinates ($x, y$) and one angular ($\theta$).
  \item A rigid body in a plane has \textbf{three DOF}.
\end{itemize}

\vspace{4pt}
\textbf{A rigid body in 3-D space:}
\begin{itemize}
  \item Needs 6 parameters: three lengths ($x, y, z$) plus three angles ($\theta, \phi, \rho$).
  \item A rigid body in three-space has \textbf{six DOF}.
\end{itemize}

\vspace{4pt}
\textbf{Key assumption:} For kinematic analysis, all bodies are \textbf{rigid} and \textbf{massless}. Deformations can be superposed later.
\pref{30--31}
\end{frame}

% ============================================================
% 2.2 TYPES OF MOTION
% ============================================================

\begin{frame}
\frametitle{Types of Motion in a Plane}

\begin{yellowbox}
\textbf{Pure rotation:} One point (center of rotation) has no motion w.r.t.\ the stationary frame. All other points describe arcs about that center. A reference line changes only its angular orientation.
\end{yellowbox}

\vspace{4pt}
\begin{yellowbox}
\textbf{Pure translation:} All points on the body describe parallel (curvilinear or rectilinear) paths. A reference line changes linear position but \emph{not} angular orientation.
\end{yellowbox}

\vspace{4pt}
\begin{yellowbox}
\textbf{Complex motion:} A simultaneous combination of rotation and translation. Any reference line changes both its linear position and angular orientation. At every instant there exists a (moving) center of rotation.
\end{yellowbox}

\vspace{4pt}
Translation and rotation are \textbf{independent} motions. In a 2-D coordinate system: $x$, $y$ = translation components; $\theta$ = rotation component.
\pref{32}
\end{frame}

% ============================================================
% 2.3 LINKS, JOINTS, KINEMATIC CHAINS
% ============================================================

\begin{frame}
\frametitle{Links}

All common mechanisms (linkages, cams, gears, belts, chains) are variations on links and joints.

\vspace{6pt}
\begin{yellowbox}
A \textbf{link} is an (assumed) rigid body that possesses at least two \textbf{nodes} --- points for attachment to other links.
\end{yellowbox}

\vspace{6pt}
Links are classified by their number of nodes (link order):

\begin{center}
\begin{tabular}{lcl}
\toprule
\textbf{Name} & \textbf{Nodes} & \textbf{Description} \\
\midrule
Binary link & 2 & Simplest; two attachment points \\
Ternary link & 3 & Three attachment points \\
Quaternary link & 4 & Four attachment points \\
\bottomrule
\end{tabular}
\end{center}

\vspace{4pt}
Real links can be any shape; a ``kinematic'' link is just a line between nodes. Shading or crosshatching indicates a solid link in kinematic diagrams.
\pref{32--33}
\end{frame}

\begin{frame}
\frametitle{Joints (Kinematic Pairs)}

\begin{yellowbox}
A \textbf{joint} is a connection between two or more links (at their nodes) that allows some motion or potential motion between them. Also called a \textbf{kinematic pair}.
\end{yellowbox}

\vspace{4pt}
Joints can be classified by:
\begin{enumerate}
  \item Type of contact (line, point, surface)
  \item Number of DOF allowed at the joint
  \item Type of physical closure (\textbf{force} or \textbf{form} closed)
  \item Number of links joined (\textbf{joint order})
\end{enumerate}

\vspace{4pt}
\textbf{Lower pairs} (Reuleaux): surface contact $\to$ better lubrication, longer life.

\textbf{Higher pairs}: point or line contact.

\vspace{4pt}
For \textbf{planar mechanisms}, only two lower pairs are usable:
\begin{itemize}
  \item \textbf{Revolute (R)} --- pin joint, 1 DOF (rotation)
  \item \textbf{Prismatic (P)} --- slider joint, 1 DOF (translation)
\end{itemize}
\pref{33--34}
\end{frame}

\begin{frame}
\frametitle{The Six Lower Pairs}

\begin{center}
\small
\begin{tabular}{l|c|c|l}
\toprule
\textbf{Name} & \textbf{Symbol} & \textbf{DOF} & \textbf{Contains} \\
\midrule
Revolute & R & 1 & R \\
Prismatic & P & 1 & P \\
Helical (screw) & H & 1 & R + P (coupled) \\
Cylindric & C & 2 & R + P \\
Spherical (ball) & S & 3 & R + R + R \\
Planar (flat) & F & 3 & R + P + P \\
\bottomrule
\end{tabular}
\end{center}

\vspace{6pt}
\textbf{R} and \textbf{P} are the basic building blocks of all other pairs.

\vspace{4pt}
\textbf{Full joints} (lower pairs): 1 DOF --- remove 2 DOF from the system.

\textbf{Half joints} (higher pairs, roll-slide): 2 DOF --- remove only 1 DOF.

\vspace{4pt}
\textbf{Joint order} $=$ number of links joined minus one. Two links $\to$ order 1. Three links $\to$ order 2 (a \textbf{multiple joint}).
\pref{33--35}
\end{frame}

\begin{frame}
\frametitle{Form Closure vs.\ Force Closure}

\begin{yellowbox}
\textbf{Form-closed joint:} Kept together by its geometry (e.g., pin in hole, slider in slot).

\textbf{Force-closed joint:} Requires an external force to keep it together (e.g., pin in half-bearing, cam-follower with spring return).
\end{yellowbox}

\vspace{6pt}
The choice matters for mechanism behavior:
\begin{itemize}
  \item In \textbf{linkages}: form closure is usually preferred (easy to accomplish).
  \item In \textbf{cam-follower systems}: force closure is often preferred (spring-loaded followers).
\end{itemize}

\vspace{6pt}
\textbf{Rolling contact} (e.g., tire on road): can be \textbf{pure roll} (full joint, 1 DOF), \textbf{pure slide} (full joint, 1 DOF), or \textbf{roll-slide} (half joint, 2 DOF), depending on friction.
\pref{35}
\end{frame}

\begin{frame}
\frametitle{Kinematic Chains, Mechanisms, and Machines}

\begin{yellowbox}
\textbf{Kinematic chain:} An assemblage of links and joints interconnected to provide a controlled output motion in response to a supplied input motion.
\end{yellowbox}

\vspace{4pt}
\begin{yellowbox}
\textbf{Mechanism:} A kinematic chain in which at least one link is \textbf{grounded} (attached to the frame of reference).
\end{yellowbox}

\vspace{4pt}
\begin{yellowbox}
\textbf{Machine} (Reuleaux): A collection of mechanisms arranged to transmit forces and do work.
\end{yellowbox}

\vspace{6pt}
Special link types:
\begin{itemize}
  \item \textbf{Crank:} makes a complete revolution; pivoted to ground.
  \item \textbf{Rocker:} oscillatory (back-and-forth) rotation; pivoted to ground.
  \item \textbf{Coupler} (connecting rod): complex motion; \emph{not} pivoted to ground.
  \item \textbf{Ground:} any link or links that are fixed (nonmoving).
\end{itemize}
\pref{35--36}
\end{frame}

% ============================================================
% 2.4 KINEMATIC DIAGRAMS
% ============================================================

\begin{frame}
\frametitle{Drawing Kinematic Diagrams}

Kinematic diagrams are simplified, schematic representations of mechanisms. They show only the kinematic structure (links and joints), not the physical shape.

\vspace{6pt}
\textbf{Notation conventions:}
\begin{itemize}
  \item Binary, ternary, quaternary links: shown as lines or shaded polygons.
  \item \textbf{Moving revolute joint:} small circle.
  \item \textbf{Grounded revolute joint:} circle with hatching/ground symbol.
  \item \textbf{Moving slider:} rectangle on a line.
  \item \textbf{Grounded slider:} rectangle with ground symbol.
  \item \textbf{Half joints:} shown with both contact surfaces.
\end{itemize}

\vspace{4pt}
\begin{redbox}
It is \textbf{critical} that your diagram clearly indicate which links/joints are grounded and which can move. Without this, nobody can interpret your design's kinematics.
\end{redbox}
\pref{36--37}
\end{frame}

% ============================================================
% 2.5 GRUEBLER'S EQUATION
% ============================================================

\begin{frame}
\frametitle{Determining DOF: Gruebler's Equation}

\begin{yellowbox}
\textbf{Degree of freedom} (mobility $M$): The number of inputs needed to provide a predictable output; equivalently, the number of independent coordinates required to define the system's position.
\end{yellowbox}

\vspace{4pt}
Kinematic chains can be \textbf{open} (like a robot arm) or \textbf{closed} (loop).

\vspace{4pt}
Any link in a plane has 3 DOF. A system of $L$ unconnected links: $3L$ DOF. Each full joint removes 2 DOF. Each grounded link removes 3 DOF.

\vspace{4pt}
\begin{redbox}
\textbf{Gruebler's equation:}
\[
  M = 3L - 2J - 3G \tag{2.1a}
\]
where $M$ = mobility, $L$ = number of links, $J$ = number of joints, $G$ = number of grounded links.
\end{redbox}

Since $G$ is always 1 in any real mechanism:
\[
  M = 3(L - 1) - 2J \tag{2.1b}
\]
\pref{38--40}
\end{frame}

\begin{frame}
\frametitle{Kutzbach's Modification}

Half joints remove only 1 DOF (they have 2 DOF). Using \textbf{Kutzbach's} form:

\begin{redbox}
\[
  M = 3(L - 1) - 2J_1 - J_2 \tag{2.1c}
\]
where:
\begin{itemize}
  \item $J_1$ = number of 1-DOF (full) joints
  \item $J_2$ = number of 2-DOF (half) joints
\end{itemize}
\end{redbox}

\vspace{6pt}
\textbf{Important counting rules:}
\begin{itemize}
  \item Half joints count as 1/2 in the $J$ of equation~2.1b.
  \item \textbf{Multiple joints} (3 or more links at one pin) count as one less than the number of links joined. E.g., 3 links at one pin $= 2$ full joints.
  \item This equation uses no information about link \textbf{sizes or shapes} --- only their \textbf{quantity}.
\end{itemize}
\pref{40}
\end{frame}

\begin{frame}
\frametitle{Example: Applying Gruebler's Equation}

\textbf{Example 1:} A linkage with $L = 8$, $J = 10$ (all full joints, including multiple joints counted correctly), no half joints.
\[
  M = 3(8-1) - 2(10) = 21 - 20 = 1 \quad \checkmark
\]
One input needed to drive the mechanism.

\vspace{6pt}
\textbf{Example 2:} A linkage with $L = 6$, $J_1 = 7$ (full), $J_2 = 1$ (half joint).
\[
  M = 3(6-1) - 2(7) - 1 = 15 - 14 - 1 = 0
\]
This is a \textbf{structure}, not a mechanism.

\vspace{6pt}
\begin{greenbox}
\textbf{Quick check:} For any proposed mechanism, compute $M$ from Gruebler's equation \emph{before} investing time in detailed design. It requires no information about link sizes, only counts.
\end{greenbox}
\pref{40--41}
\end{frame}

% ============================================================
% 2.6 MECHANISMS AND STRUCTURES
% ============================================================

\begin{frame}
\frametitle{Mechanisms, Structures, and Preloaded Structures}

The DOF completely predicts the character of a link assembly:

\begin{center}
\begin{tabular}{c|l|l}
\toprule
\textbf{DOF} & \textbf{Type} & \textbf{Behavior} \\
\midrule
$M > 0$ & Mechanism & Links have relative motion \\
$M = 0$ & Structure & No motion possible (rigid) \\
$M < 0$ & Preloaded structure & No motion; internal stresses \\
\bottomrule
\end{tabular}
\end{center}

\vspace{6pt}
\begin{itemize}
  \item $M > 0$: a mechanism with $M$ inputs needed.
  \item $M = 0$: \textbf{exact constraint} --- a structure (like a truss). Parts can be assembled without strain if manufactured correctly.
  \item $M < 0$: \textbf{preloaded/indeterminate} structure --- parts must be forced together, creating internal stresses. Common in buildings, bridges, and machine frames.
\end{itemize}

\vspace{4pt}
Our concern: devices with \textbf{positive DOF} (mechanisms).
\pref{41--42}
\end{frame}

% ============================================================
% 2.7 NUMBER SYNTHESIS
% ============================================================

\begin{frame}
\frametitle{Number Synthesis}

\begin{yellowbox}
\textbf{Number synthesis:} The determination of the number and order of links and joints necessary to produce motion of a particular DOF.
\end{yellowbox}

\vspace{4pt}
\textbf{Key result:} For full joints only, an odd number of DOF requires an \emph{even} number of links, and vice versa.

\vspace{4pt}
From Gruebler: for 1-DOF mechanisms with only full joints, $L$ must be even: $L = 4, 6, 8, \ldots$

\vspace{6pt}
Important relationships for link orders ($B$ = binary, $T$ = ternary, $Q$ = quaternary, \ldots):
\[
  L = B + T + Q + P + H + \cdots \tag{2.4a}
\]
\[
  M = B - Q - 2P - 3H - 3 \tag{2.4e}
\]

\begin{redbox}
The ternary links \textbf{drop out} of the mobility equation! Adding or subtracting ternary links \emph{in pairs} does not affect the DOF of the mechanism.
\end{redbox}
\pref{42--44}
\end{frame}

\begin{frame}
\frametitle{Number Synthesis: Results for 1-DOF Mechanisms}

For 1-DOF ($M=1$) using only full revolute joints, up to 8 links:

\begin{center}
\small
\begin{tabular}{c|ccccc}
\toprule
$L$ & Binary & Ternary & Quaternary & Pentagonal & Hexagonal \\
\midrule
4 & 4 & 0 & 0 & 0 & 0 \\
\midrule
6 & 4 & 2 & 0 & 0 & 0 \\
6 & 5 & 0 & 1 & 0 & 0 \\
\midrule
8 & 7 & 0 & 0 & 0 & 1 \\
8 & 4 & 4 & 0 & 0 & 0 \\
8 & 5 & 2 & 1 & 0 & 0 \\
8 & 6 & 0 & 2 & 0 & 0 \\
8 & 6 & 1 & 0 & 1 & 0 \\
\bottomrule
\end{tabular}
\end{center}

\vspace{4pt}
\begin{yellowbox}
\textbf{Key takeaway:} The simplest 1-DOF mechanism is the \textbf{fourbar linkage} (4 binary links). For 6 links there are 2 configurations; for 8 links there are 5.
\end{yellowbox}
\pref{44--46; Table 2-2}
\end{frame}

% ============================================================
% 2.8 PARADOXES
% ============================================================

\begin{frame}
\frametitle{Gruebler Paradoxes}

Because Gruebler's criterion ignores link sizes and shapes, it can give \textbf{misleading results} for certain special geometric configurations.

\vspace{6pt}
\textbf{Example 1: The E-quintet} (5 links, 6 joints, $M = 0$ by Gruebler). But with equispaced nodes and equal-length ternary links, it actually moves ($M = 1$)!

\vspace{4pt}
\textbf{Example 2: Rolling cylinders} (3 links, 3 full joints if no slip, $M = 0$ by Gruebler). But if the center distance equals the sum of radii, it actually moves ($M = 1$).

\vspace{6pt}
\begin{redbox}
\textbf{Warning:} None of the simple mobility equations (Gruebler, Kutzbach) can resolve all paradoxes. A complete \textbf{position analysis} (Chapter~4) is necessary to guarantee mobility. The designer must be alert to these inconsistencies.
\end{redbox}
\pref{46--47}
\end{frame}

% ============================================================
% 2.9 ISOMERS
% ============================================================

\begin{frame}
\frametitle{Isomers}

\begin{yellowbox}
\textbf{Isomer} (from Greek: ``having equal parts''): Linkages that have the same number and type of links but are interconnected differently, yielding different motion properties.
\end{yellowbox}

\vspace{4pt}
Analogy: chemical isomers (n-butane vs.\ isobutane: same atoms, different connections, different properties).

\vspace{6pt}
\begin{center}
\small
\begin{tabular}{c|c}
\toprule
\textbf{Links} & \textbf{Valid isomers} \\
\midrule
4 & 1 \\
6 & 2 (Watt's chain, Stephenson's chain) \\
8 & 16 \\
10 & 230 \\
12 & 6856 \\
\bottomrule
\end{tabular}
\end{center}

\vspace{4pt}
An isomer is \textbf{invalid} if the highest-order link has more than $n/2$ nodes (where $n$ = total links), or if it contains a structural subchain with $DOF = 0$.
\pref{47--49; Table 2-3}
\end{frame}

% ============================================================
% 2.10 LINKAGE TRANSFORMATION
% ============================================================

\begin{frame}
\frametitle{Linkage Transformation Rules}

Given a basic linkage from number synthesis (all revolute joints), we can transform it into a wider variety of mechanisms:

\begin{yellowbox}
\textbf{Six transformation rules:}
\begin{enumerate}
  \item Revolute joints $\to$ prismatic joints (if $\geq 2$ revolutes remain in the loop). \textbf{DOF unchanged.}
  \item Full joint $\to$ half joint. \textbf{DOF increases by one.}
  \item Remove a link. \textbf{DOF decreases by one.}
  \item Rules 2 + 3 combined. \textbf{DOF unchanged.}
  \item Partially shrink a higher-order link (coalesce nodes $\to$ multiple joint). \textbf{DOF unchanged.}
  \item Completely shrink a higher-order link (remove it). \textbf{DOF decreases.}
\end{enumerate}
\end{yellowbox}
\pref{49--52}
\end{frame}

\begin{frame}
\frametitle{Linkage Transformation: Key Examples}

\textbf{Rule 1:} Fourbar crank-rocker $\to$ \textbf{crank-slider}

Replace the rocker pivot with a slider. The slider block is equivalent to an infinitely long rocker pivoted at infinity. Still 1 DOF. Used in piston engines and pumps.

\vspace{6pt}
\textbf{Rule 4:} Fourbar $\to$ \textbf{cam-follower}

Remove link~3, replace the full joint with a half joint between links 2 and 4. The coupler becomes an effective link of variable length. Still 1 DOF. The cam-follower is a fourbar in disguise!

\vspace{6pt}
\textbf{Rule 4:} Fourbar slider-crank $\to$ \textbf{Scotch yoke}

Replace coupler with a half joint (slot). Gives exact simple harmonic motion of the slider for constant-speed crank input.

\vspace{4pt}
\begin{greenbox}
\textbf{Insight:} All common mechanism types (linkages, cams, gears) are transformations of a common set of basic kinematic chains.
\end{greenbox}
\pref{50--52}
\end{frame}

% ============================================================
% 2.11 INTERMITTENT MOTION
% ============================================================

\begin{frame}
\frametitle{Intermittent Motion}

\begin{yellowbox}
\textbf{Intermittent motion:} A sequence of motions and \textbf{dwells} (periods where the output is stationary while the input continues to move).
\end{yellowbox}

\vspace{6pt}
\textbf{Geneva mechanism:}
\begin{itemize}
  \item A transformed fourbar (coupler replaced by half joint).
  \item Input crank with a pin enters radial slots on the Geneva wheel.
  \item Produces intermittent rotation of the output.
  \item Number of slots $=$ number of ``stops'' (minimum 3).
  \item Arc on crank engages periphery to lock the wheel between stops.
\end{itemize}

\vspace{4pt}
\textbf{Ratchet and pawl:}
\begin{itemize}
  \item Driving pawl indexes the ratchet wheel (counterclockwise only).
  \item Locking pawl prevents reverse rotation.
  \item Widely used: wrenches, winches, etc.
\end{itemize}
\pref{53--54}
\end{frame}

% ============================================================
% 2.12 INVERSION
% ============================================================

\begin{frame}
\frametitle{Inversion}

\begin{yellowbox}
An \textbf{inversion} is created by grounding a \emph{different link} in the kinematic chain. A linkage has as many inversions as it has links.
\end{yellowbox}

\vspace{4pt}
Not all inversions produce distinct motions. We call those that do \textbf{distinct inversions}.

\vspace{6pt}
\textbf{The fourbar slider-crank} has 4 inversions (all distinct):

\begin{center}
\small
\begin{tabular}{c|l|l}
\toprule
\textbf{\#} & \textbf{Grounded link} & \textbf{Application} \\
\midrule
1 & Link 1 (frame) & Piston engines, pumps \\
2 & Link 2 (crank) & Whitworth/crank-shaper \\
3 & Link 3 (coupler) & Slider block rotates \\
4 & Link 4 (slider) & Hand-operated well pump \\
\bottomrule
\end{tabular}
\end{center}

\vspace{4pt}
\textbf{Watt's sixbar:} 2 distinct inversions. \textbf{Stephenson's sixbar:} 3 distinct inversions. \textbf{Pin-jointed fourbar:} 4 distinct inversions (crank-rocker, double-crank, double-rocker, triple-rocker).
\pref{53--56}
\end{frame}

% ============================================================
% 2.13 THE GRASHOF CONDITION
% ============================================================

\begin{frame}
\frametitle{The Grashof Condition}

The fourbar linkage is the simplest pin-jointed mechanism for single-DOF controlled motion. The \textbf{Grashof condition} predicts its rotation behavior based on link lengths alone.

\vspace{4pt}
Let: $S$ = shortest link, $L$ = longest link, $P$, $Q$ = other two links.

\begin{redbox}
\textbf{Grashof condition:}
\[
  S + L \leq P + Q \tag{2.8}
\]
If this inequality holds, the linkage is \textbf{Grashof} and at least one link can make a full revolution (\textbf{Class I}).

If $S + L > P + Q$: \textbf{non-Grashof} (\textbf{Class II}) --- no link can fully rotate.

If $S + L = P + Q$: \textbf{special-case Grashof} (\textbf{Class III}) --- change points exist.
\end{redbox}

\vspace{4pt}
The Grashof condition is checked on \textbf{unassembled links} --- independent of assembly order.
\pref{56}
\end{frame}

\begin{frame}
\frametitle{Grashof Class I: Inversions}

For Class I ($S + L < P + Q$), the inversions depend on which link is grounded (relative to the shortest link $S$):

\vspace{6pt}
\begin{center}
\small
\begin{tabular}{l|l|l}
\toprule
\textbf{Grounded link} & \textbf{Type} & \textbf{Code} \\
\midrule
Link adjacent to $S$ & \textbf{Crank-rocker} & GCRR \\
Shortest link $S$ & \textbf{Double-crank} (drag link) & GCCC \\
Link opposite $S$ & \textbf{Double-rocker} (coupler rotates) & GRCR \\
\bottomrule
\end{tabular}
\end{center}

\vspace{6pt}
\textbf{Class II} ($S + L > P + Q$): All four inversions are \textbf{triple-rockers} (RRR1--RRR4). No link can fully rotate.

\vspace{6pt}
\textbf{Class III} ($S + L = P + Q$): \textbf{Change points} occur when all links become colinear. Behavior is unpredictable at these points (``uncertainty configurations''). Includes the \textbf{parallelogram} and \textbf{deltoid} (kite) forms.
\pref{57--59}
\end{frame}

\begin{frame}
\frametitle{Special-Case Grashof Linkages}

\textbf{Parallelogram linkage} ($S + L = P + Q$, opposite links equal):
\begin{itemize}
  \item Coupler is in curvilinear translation (all points trace identical circles).
  \item Used to couple two output rockers (e.g., windshield wipers).
  \item Has change points every $180^\circ$.
\end{itemize}

\vspace{4pt}
\textbf{Antiparallelogram} (``butterfly'' or ``bow-tie''):
\begin{itemize}
  \item Also a double-crank but with different angular velocity at output.
  \item Change points require a companion linkage $90^\circ$ out of phase.
\end{itemize}

\vspace{4pt}
\textbf{Deltoid} (kite):
\begin{itemize}
  \item Two pairs of adjacent equal links.
  \item Short crank makes 2 revolutions per 1 revolution of the long crank.
  \item Also called an isoceles linkage or \textbf{Galloway} mechanism.
\end{itemize}

\vspace{4pt}
There is nothing inherently bad about any Grashof class --- all are useful in the right application.
\pref{59}
\end{frame}

\begin{frame}
\frametitle{Barker's Complete Classification (14 Types)}

Barker classifies all fourbar linkages using \textbf{link ratios} $\lambda_1 = r_1/r_2$, $\lambda_3 = r_3/r_2$, $\lambda_4 = r_4/r_2$ (where $r_1$ = ground, $r_2$ = input, $r_3$ = coupler, $r_4$ = output):

\vspace{4pt}
\begin{center}
\scriptsize
\begin{tabular}{c|c|l|l}
\toprule
\textbf{Type} & \textbf{Class} & \textbf{Barker's Designation} & \textbf{Common Name} \\
\midrule
1--4 & I (Grashof) & GCCC, GCRR, GRCR, GRRC & Double-crank, crank-rocker, etc. \\
5--8 & II (non-Grashof) & RRR1--RRR4 & All triple-rockers \\
9--12 & III (special) & SCCC, SCRR, SRCR, SRRC & Change-point types \\
13 & III & S2X & Parallelogram / deltoid \\
14 & III & S3X & Square \\
\bottomrule
\end{tabular}
\end{center}

\vspace{4pt}
The solution space ($\lambda_1, \lambda_3, \lambda_4$) is divided into 8 volumes by 4 zero-mobility planes and 3 change-point planes. Each volume contains a unique Barker type.
\pref{60--61; Table 2-4}
\end{frame}

% ============================================================
% 2.14 LINKAGES > 4 BARS
% ============================================================

\begin{frame}
\frametitle{Linkages of More Than Four Bars}

\textbf{Geared fivebar linkage:}
\begin{itemize}
  \item Add one link and one joint to a fourbar $\to$ 2 DOF.
  \item Add a gearset between two links $\to$ back to 1 DOF.
  \item Called the \textbf{geared fivebar mechanism (GFBM)}.
  \item More complex motions than the fourbar.
\end{itemize}

\vspace{6pt}
\textbf{Sixbar linkages} --- two fundamental types:
\begin{itemize}
  \item \textbf{Watt's sixbar:} two fourbars connected \emph{in series} (sharing two links).
  \item \textbf{Stephenson's sixbar:} two fourbars connected \emph{in parallel} (sharing two links).
\end{itemize}

\vspace{4pt}
Many linkages can be designed by combining multiple fourbar chains as \textbf{basic building blocks}. The basic building block of zero-DOF structures is the \textbf{delta triplet} (3-link triangle).

\vspace{4pt}
\textbf{Grashof-type criteria} exist for geared fivebars (Ting):
\[
  L_1 + L_2 + L_5 < L_3 + L_4 \tag{2.12}
\]
\pref{62--64}
\end{frame}

\begin{frame}
\frametitle{N-Bar Rotatability Criteria}

For general $N$-bar revolute-jointed linkages (Ting et al.):

\vspace{4pt}
\textbf{Assemblability} (necessary and sufficient): the longest link must be shorter than the sum of all others:
\[
  L_N \leq \sum_{k=1}^{N-1} L_k \tag{2.13}
\]

\textbf{Rotatability classification} for single-loop, $N$-bar chains:
\begin{yellowbox}
\begin{itemize}
  \item \textbf{Class I:} $L_N + (L_1 + L_2 + \cdots + L_{N-3}) < L_{N-2} + L_{N-1}$
  \item \textbf{Class II:} $>$ (non-Grashof)
  \item \textbf{Class III:} $=$ (change points)
\end{itemize}
\end{yellowbox}

\vspace{4pt}
\textbf{Revolvability} of any link $L_i$:
\[
  L_i + L_N \leq \sum_{\substack{k=1 \\ k \neq i}}^{N-1} L_k \tag{2.16}
\]
If $L_i$ is revolvable, any link not longer than $L_i$ is also revolvable.
\pref{64}
\end{frame}

% ============================================================
% 2.15--2.17 SPRINGS, COMPLIANT, MEMS
% ============================================================

\begin{frame}
\frametitle{Springs as Links, Compliant Mechanisms, and MEMS}

\textbf{Springs as links:}
\begin{itemize}
  \item A spring acts as an additional link of variable length.
  \item Can counterbalance static loads to reduce effective DOF to zero, holding the mechanism in static equilibrium.
  \item Example: desk lamp, automobile hood strut.
  \item Spring constant: $k = F/x$.
\end{itemize}

\vspace{4pt}
\textbf{Compliant mechanisms:}
\begin{itemize}
  \item Use \textbf{flexibility} of members instead of discrete joints.
  \item Fewer parts, no assembly, inherent spring loading.
  \item Examples: plastic box ``living hinge,'' compliant forceps, toggle switches.
  \item Difficult to analyze (large deflections).
\end{itemize}

\vspace{4pt}
\textbf{MEMS} (Micro Electro-Mechanical Systems):
\begin{itemize}
  \item Mechanisms at the micrometer scale, etched from silicon wafers.
  \item Compliant mechanisms are well-suited to MEMS fabrication.
  \item Applications: airbag sensors, pressure monitors, microswitches.
\end{itemize}
\pref{65--68}
\end{frame}

% ============================================================
% 2.18 PRACTICAL CONSIDERATIONS
% ============================================================

\begin{frame}
\frametitle{Practical Considerations: Joint Types}

\textbf{Revolute (pin) joints:}
\begin{itemize}
  \item Preferred for most applications: easy to design, lubricate, and manufacture.
  \item Use a \textbf{journal} (sleeve) bearing for simplicity, or ball/roller bearings for precision.
  \item \textbf{Spherical rod ends}: accommodate misalignment; allow link length adjustment.
  \item \textbf{Pillow blocks} and \textbf{flange mounts}: convenient for grounded pivots.
\end{itemize}

\vspace{6pt}
\textbf{Prismatic (slider) joints:}
\begin{itemize}
  \item Require machined slots or rods; lubrication is more difficult.
  \item Use linear ball bearings for precision applications.
\end{itemize}

\vspace{6pt}
\textbf{Higher (half) joints:}
\begin{itemize}
  \item Cam-follower and pin-in-slot: line contact, worse lubrication.
  \item Must be run in oil bath with seals for long life.
\end{itemize}
\pref{69--71}
\end{frame}

\begin{frame}
\frametitle{Practical Considerations: Mounting and Bearing Ratio}

\textbf{Cantilever vs.\ straddle mount:}
\begin{itemize}
  \item \textbf{Straddle mount} (pin supported on both sides): preferred --- pin is in double shear, no bending moment on links.
  \item \textbf{Cantilever mount} (pin supported on one side only): weaker --- single shear, bending moment. Use a shoulder screw if unavoidable.
\end{itemize}

\vspace{6pt}
\begin{redbox}
\textbf{Bearing ratio} for slider joints:
\[
  BR = \frac{L}{D} = \frac{\text{effective length of slider}}{\text{effective diameter of bearing}}
\]
For smooth operation: $BR \geq 1.5$. The larger the better.

A drawer with only side-rails as guides has $BR < 1$ and will jam. Adding a center guide dramatically increases the effective $L/D$.
\end{redbox}
\pref{71--73}
\end{frame}

\begin{frame}
\frametitle{Linkages vs.\ Cams: Design Trade-offs}

\begin{center}
\small
\begin{tabular}{l|c|c}
\toprule
\textbf{Criterion} & \textbf{Linkages} & \textbf{Cams} \\
\midrule
Size relative to output & Larger & Smaller \\
Synthesis difficulty & Hard & Easy (with software) \\
Manufacturing cost & Low & High (precision grinding) \\
Precision & Lower & Higher \\
Dwell capability & Difficult & Easy \\
Environmental tolerance & Excellent & Poor (needs sealing) \\
Lubrication needs & Minimal & Critical \\
Speed capability & Higher & Lower \\
Load capacity & Higher & Lower \\
Sensitivity to mfg.\ error & Low & High \\
\bottomrule
\end{tabular}
\end{center}

\vspace{4pt}
\begin{yellowbox}
\textbf{Design rule:} Always try a pure linkage solution first. Use cams only when a straight, sliding, or exact dwell motion is required that a linkage cannot provide economically.
\end{yellowbox}
\pref{73}
\end{frame}

% ============================================================
% 2.19 MOTORS AND DRIVERS
% ============================================================

\begin{frame}
\frametitle{Motors and Drivers: Overview}

Unless manually operated, a mechanism needs a \textbf{driver} (motor, engine, cylinder). The main division:

\vspace{4pt}
\begin{center}
\begin{tabular}{l|l}
\toprule
\textbf{Rotary input} & \textbf{Linear input} \\
\midrule
Electric motors & Hydraulic/pneumatic cylinders \\
Gasoline/diesel engines & Solenoids \\
\bottomrule
\end{tabular}
\end{center}

\vspace{6pt}
\textbf{Electric motors} --- classified by:
\begin{itemize}
  \item Function: \textbf{gearmotors}, \textbf{servomotors}, \textbf{stepper motors}
  \item Configuration: \textbf{AC} vs.\ \textbf{DC}
  \item Power class: subfractional ($<$1/20 HP), fractional (1/20--1 HP), integral ($>$1 HP)
\end{itemize}

\vspace{4pt}
A motor is selected based on its \textbf{torque-speed curve}: this predicts how speed varies with applied load.
\pref{74--75}
\end{frame}

\begin{frame}
\frametitle{DC Motor Types}

All DC motors share: max torque at stall (zero speed), zero torque at max speed. Power = torque $\times$ angular velocity.

\vspace{6pt}
\begin{center}
\small
\begin{tabular}{l|l|l}
\toprule
\textbf{Type} & \textbf{Characteristic} & \textbf{Application} \\
\midrule
Permanent magnet & Torque varies linearly with speed & General purpose \\
Shunt-wound & Flatter curve, speed-sensitive & Fans, blowers \\
Series-wound & Very high start torque, can run away & Traction drives \\
Compound-wound & Combines shunt + series & Cranes, hoists \\
\bottomrule
\end{tabular}
\end{center}

\vspace{4pt}
\textbf{Speed-controlled DC motors:} Use a controller to maintain constant speed despite varying load. Can also use a \textbf{flywheel} to smooth speed variations.

\vspace{4pt}
\textbf{AC motors:} Cheaper for continuous rotation. Synchronous speed: $n_s = 120f / p$ (where $f$ = line frequency, $p$ = poles). Most common: 4-pole $\to$ 1725 rpm (async) at 60 Hz.
\pref{75--77}
\end{frame}

\begin{frame}
\frametitle{Servomotors, Steppers, and Linear Actuators}

\textbf{Servomotors:}
\begin{itemize}
  \item \textbf{Closed-loop} control (sensors feed back position/velocity).
  \item High torque, flat speed-torque curve, ``soft start'' capability.
  \item Used in: CNC machines, robots, aircraft control surfaces.
  \item Available in AC and DC types.
\end{itemize}

\vspace{4pt}
\textbf{Stepper motors:}
\begin{itemize}
  \item \textbf{Open-loop} (no feedback --- can lose synchronization).
  \item Intermittent motion: one step per pulse (200--2000+ steps/rev).
  \item High holding torque but low drive torque.
  \item Good for positioning (printers, small actuators).
\end{itemize}

\vspace{4pt}
\textbf{Linear actuators:}
\begin{itemize}
  \item \textbf{Hydraulic/pneumatic cylinders:} limited stroke, straight-line output. Common in construction equipment (become links in a slider-crank).
  \item \textbf{Solenoids:} very short stroke (2--3 cm), fast response, used as switches.
\end{itemize}
\pref{78--80}
\end{frame}

% ============================================================
% SUMMARY
% ============================================================

\begin{frame}
\frametitle{Chapter 2: Summary of Key Concepts}

\begin{enumerate}
  \item A rigid body in a plane has \textbf{3 DOF}; in 3-D space, \textbf{6 DOF}.
  \item \textbf{Links} (binary, ternary, quaternary) connect via \textbf{joints} (revolute, prismatic, half joints).
  \item \textbf{Gruebler/Kutzbach}: $M = 3(L-1) - 2J_1 - J_2$. Positive $M$ = mechanism; zero = structure; negative = preloaded.
  \item \textbf{Number synthesis}: only even numbers of links for odd-DOF mechanisms. Simplest 1-DOF = fourbar (4 binary links).
  \item \textbf{Linkage transformation} (6 rules) converts basic chains into sliders, cams, gears, etc.
  \item \textbf{Grashof condition}: $S + L \leq P + Q$. Determines whether a full revolution is possible. 3 classes, 14 Barker types.
  \item \textbf{Inversion}: grounding different links creates different motions from the same chain.
  \item \textbf{Motor selection}: based on torque-speed characteristics. Servos (closed-loop) vs.\ steppers (open-loop).
\end{enumerate}
\end{frame}

\begin{frame}
\frametitle{Key Equations to Remember}

\begin{yellowbox}
\begin{tabular}{@{}ll@{}}
\textbf{Gruebler's eq.} & $M = 3(L - 1) - 2J$ \\[4pt]
\textbf{Kutzbach's mod.} & $M = 3(L - 1) - 2J_1 - J_2$ \\[4pt]
\textbf{Grashof condition} & $S + L \leq P + Q$ \\[4pt]
\textbf{Spatial mobility} & $M = 6(L-1) - 5J_1 - 4J_2 - 3J_3 - 2J_4 - J_5$ \\[4pt]
\textbf{Assemblability ($N$-bar)} & $L_N \leq \sum_{k=1}^{N-1} L_k$ \\[4pt]
\textbf{Revolvability} & $L_i + L_N \leq \sum_{k=1, k\neq i}^{N-1} L_k$ \\[4pt]
\textbf{AC motor speed} & $n_s = 120f / p$ \\[4pt]
\textbf{Bearing ratio} & $BR = L/D \geq 1.5$ \\
\end{tabular}
\end{yellowbox}
\end{frame}

\begin{frame}
\frametitle{Key Definitions to Remember}

\begin{yellowbox}
\small
\begin{tabular}{@{}p{3cm}p{7.5cm}@{}}
\textbf{DOF / Mobility} & Number of independent inputs needed \\
\textbf{Link} & Rigid body with $\geq 2$ nodes \\
\textbf{Joint (pair)} & Connection allowing relative motion \\
\textbf{Full joint} & 1 DOF (removes 2 DOF from system) \\
\textbf{Half joint} & 2 DOF (removes 1 DOF from system) \\
\textbf{Kinematic chain} & Assemblage of links and joints \\
\textbf{Mechanism} & Grounded kinematic chain \\
\textbf{Crank} & Full revolution, pivoted to ground \\
\textbf{Rocker} & Oscillation, pivoted to ground \\
\textbf{Coupler} & Complex motion, not grounded \\
\textbf{Inversion} & Grounding a different link \\
\textbf{Grashof} & $S + L \leq P + Q$; full rotation possible \\
\textbf{Isomer} & Same link counts, different connections \\
\textbf{Servomotor} & Closed-loop motor with feedback \\
\end{tabular}
\end{yellowbox}
\end{frame}

\begin{frame}
\frametitle{Self-Check Questions}

Test your understanding of Chapter~2:

\begin{enumerate}
  \item How many DOF does a rigid body have in a plane? In 3-D space?
  \item What is the difference between a full joint and a half joint? How does each affect the mobility equation?
  \item State Gruebler's equation. What are $L$, $J$, and $G$?
  \item How do you count multiple joints in Gruebler's equation?
  \item A mechanism has 6 links and 7 full joints. What is its DOF?
  \item State the Grashof condition. What are the three classes?
  \item For a Grashof Class I fourbar, what type of mechanism results from grounding the shortest link? A link adjacent to it?
  \item Name the six linkage transformation rules. Which ones preserve DOF?
  \item What is the difference between Watt's sixbar and Stephenson's sixbar?
  \item What is the bearing ratio and why must it be $\geq 1.5$?
  \item Compare servomotors and stepper motors. When would you use each?
  \item What is a Gruebler paradox? Give an example.
\end{enumerate}
\end{frame}

\end{document}
